\chapter{Zusammenfassung}
\label{cha:zusammenfassung}

\section{Zusammenfassung}
\section{Mögliche Erweiterungen\textcolor{brown}{FAIK}}
Die modulare Architektur der Motorsteuerung erlaubt prinzipiell eine nahezu unbegrenzte Erweiterung des Funktionsumfangs. Im Hinblick auf die technische Realisierbarkeit in der nahen Zukunft stehen jedoch folgende Optionen im Fokus, um den Einsatzbereich des Roboters sinnvoll zu vergrößern:

\begin{itemize}
	\item \textbf{Aktiver Notbremsassistent:} Durch die Integration von Umfeldsensorik (z. B. Ultraschall) kann eine Sicherheitsfunktion implementiert werden, die den Roboter bei Hinderniserkennung automatisch stoppt, um Kollisionen zu verhindern.
	
	\item \textbf{Mensch-Maschine-Interaktion mittels lokaler KI:} Durch den Einsatz kompakter Sprachmodelle (Small Language Models), die lokal auf dem Bordcomputer ausgeführt werden, kann eine natürliche Interaktion mit dem Menschen ermöglicht werden. Dies umfasst das Verstehen von Sprachbefehlen sowie die Generierung von Antworten, ohne auf eine permanente Internetverbindung angewiesen zu sein.
	
	\item \textbf{Autonomes Fahren mittels Kameratechnik:} Die Anbindung eines Kamerasystems ermöglicht die Erweiterung um bildverarbeitende Algorithmen. Damit können Funktionen wie Spurverfolgung (Line Tracking) oder visuelle Navigation realisiert werden, um autonom von einem Start- zu einem Zielpunkt zu navigieren.
	
	\item \textbf{Remote-Steuerung und Telemetrie:} Alternativ zur lokalen Signalverarbeitung kann eine Schnittstelle für drahtlose Kommunikation (WiFi/Bluetooth) geschaffen werden. Dies ermöglicht eine Fernsteuerung über mobile Endgeräte sowie die Rückübertragung von Telemetriedaten (z. B. Akkustand, Motorauslastung) an den Bediener.
	
\end{itemize}

\section{Fazit}



\chapter{NICHT FERGESSEN}

1) IMMER ERKLÄREN WARUM WIR SO ENTSCHIEDEN HABEN UND NICHT ANDERS(hat max mehrmals betont)








FAIK:  flowchart \ref{fig:flowchart_steuerung} kleingeschriebene block-> alles einheitlich!



%
%Auf zwei bis drei Seiten soll auf folgende Punkte eingegangen werden:
%
%\begin{itemize}
%	\item Welches Ziel sollte erreicht werden
%	\item Welches Vorgehen wurde gewählt
%	\item Was wurde erreicht, zentrale Ergebnisse nennen, am besten quantitative Angaben machen
%	\item Konnten die Ergebnisse nach kritischer Bewertung zum Erreichen des Ziels oder zur Problemlösung beitragen
%	\item  Ausblick
%\end{itemize}
%
%In der Zusammenfassung sind unbedingt klare Aussagen zum Ergebnis der Arbeit zu nennen. Üblicherweise können Ergebnisse nicht nur qualitativ, sondern auch quantitativ benannt werden, z.~B. \glqq \ldots konnte eine Effizienzsteigerung von \SI{12}{\percent} erreicht werden.\grqq~oder \glqq \ldots konnte die Prüfdauer um \SI{2}{\hour} verkürzt werden\grqq.
%
%Die Ergebnisse in der Zusammenfassung sollten selbstverständlich einen Bezug zu den in der Einleitung aufgeführten Fragestellungen und Zielen haben.
